\quellseite{665}
\chapter{Index of Subjects}

Abstract 210, 250

Abstraction as an element of method 448 f., 484 f., 566 ff.

Acromatic method, defined 450, why required 450 f., when to be employed 451 f.

Action, defined 236; teaching should lead to action 261, 446, see Habit and Will.

Adolescence, age of 263 ff., characteristics 264 ff., qualifications for leadership of adolescents 268 ff.

Adult Bible Class 177, 191, see Bible Class.

Affections 230.

Agencies of education in the Church 388 ff.

Aim of religious instruction 273 ff., not Christian maturity 274, but congregational maturity 275 ff., 615 f.; precise definition of aim 280; importance of precise definition 273; references to the aim 238, 253, 373, 438, 446, 484, 528, 615 f., 630 ff.; on the history of the aim 12 f., 46, 78, 89 f., 106, 126, 275.

Aim, statement of, as part of the method of instruction 448, 467 f., 549 ff., 563, 574; the “formal aim” 468; see also Method.

America, history of Christian education in 153 ff., bibliography 153 f.; the Sunday School movement in America 165 f.; Lutheran catechetical labors 180 ff., bibliography 153 f., 182, 184 f.

Analogy 216.

Analytical development defined 442, advantages 443.

Analyzing method 135, 137, 140, 503.

Announcement for communion, instruction concerning 381.

Answers, the treatment of 221, 461.

Anthropomorphical view of God 262.

Apologetic element in instruction, the, 110, 432, 434, 445.

Apologies, \emph{Justin’s} 15, 19, 20, 22, 24.

\emph{Apostolic Constitutions} 27, 29, 35, 47, 53, 56.

Apperception 212 f., 250, 449, 480, 526.

Application as an element of method 134, 176, 253, 445 ff., 467; necessity 445 f., 486; nature 486; only one application to be made 487; examples 562, 589 f., 592, 603 f., 605 f.; bibliography of helps for application 489.

Arcane Discipline 31, 39 ff.

Archaic words 440.

Argument as an element of method 444.

Argument from prophecy and fulfilment 20, 28, 34.

Ascetic influence upon the candidate for baptism 36.

Association as an element of method 448, 467, 484 f.

Association of concepts 206 f., 440 f.; in infancy 244, in childhood 250 f., 255.

\quellseite{666}
Athanasian Creed 66.

Attention 213; in infancy 243, in childhood 250.

Augustana 121, 278.

Ave Maria as material for instruction 69, 74 ff.; Luther’s sermon on mariolatry 86.

Baptism, evaluation of, 12, 16, 21 f., 36, 40, 48, 52, 62, 275, 408, 427 f., 616; rite of, 18, 23 f., 39 ff., 47 f., 54, 63 f.; seasons for, 30, 36, 40 ff., 46, 53 f., 64, 66; place of, 47; ignored in certain lesson plans 186, 410; necessitates instruction 4, 391, 548; infant baptism 51 ff., 64, 275; instruction concerning baptism 21, 29, 37, 49, 78, 88 f., 92 ff., 106, 110, 315 f., 381; elements of baptismal instruction in the early church 18 ff., 23, 28 ff., 34 ff., 37.

Baptists 168, 172, 180.

\emph{Barnabas, the Epistle of}, 24.

Beginners’ Classes in Sunday school, 293, 391, 418; selection of material for such classes 428; bibliography of teaching helps 465 f.

\emph{Benedicite and Gratias} as material for instruction 76, 92, 94.

Bible, evaluation 368 f.; necessity of instruction concerning the Bible 276, 287, 368 f.; nature of such instruction 367 ff.; method 465 ff. and see Method. Bible reading 27, 29, 37, 42, 55 f., 67, 130, 134, 137, 147, 149, 166, 187, 371 ff., 394 ff., 416, 493, 532, 536; special periods for Bible reading 371; how to connect with other elements of religious instruction 371 f., 416, 532 f., 536; bibliography of Bible reading 367 f. School Bibles 122, 147 f., 368, 373. The Bible in the Public School 192 ff., 394 ff. Bible Biography 149, 187, 296 f., 420, 429 f., 433 f., 446, 493, 580; Bible Geography 187, 419, 430, 493, 580; Introduction into the Bible 28, 137, 149, 187, 191, 308, 372, 416, 428 f., 431 f., 434; bibliography of introduction 367 f. Bible Class 171, 177, 191, 372, 406, 431 ff.

Bible Passages as material for instruction 26, 28, 56, 89, 122 f., 129, 131, 134, 149, 157, 160, 171, 176, 182, 369, 411 ff., esp. 414, 499; bibliography 367 f. Collections of Bible passages 28, 122 f., 129, 131, 134, 182, 369.

\emph{Biblia Historica and Biblia Pauperum} 60, 68, 70, 76.

Biblical History as material for instruction 289 ff.; its necessity 286, 291; its place in the curriculum as compared to the Catechism 290 f., 308; to be given in every grade 291, 308; its purpose 295 ff., 446, 466 f.; method of teaching 464 ff.; helps for teaching esp. 465 f.; modernization of Biblical History 476; principles governing the selection of stories 292, 375, 410 ff., 419, 422, 426 ff.; the four elements contained in every story 473 f. The first Biblical History for the Christian home (Luther’s Passionale) 82, 97; the first English Biblical History (Watt’s Scripture History) 161. History of instruction in Biblical History 20 f., 28, 30, 34, 56, 64, 67 f., 76, 90, 97, 121, 129 f., 135, 137 f., 147 ff., 161, 176 ff., 181, 185, 187, 191, 193, 290, 371.

Biblical pictures and picture books 76 f., 95, 97, 114, 130, 244, 246, 252, 260, 494 ff.; bibliography 495 f.

\quellseite{667}
Biographical element in training 250, 256, 260, 384 f., see also Bible Biography.

\emph{Bishop’s Book} 155.

\emph{Book of Common Prayer} 156, 158.

\emph{Book of Concord} (1580) 127, and often.

Booklet: Baptismal 94; Catechumen’s 136, 147; Communicants’ 136; Confessional 72, 75; Laymen’s and Children’s 88; Marriage 94; Prayer 85 ff., 89, 96.

Brethren, Church of 159.

Brethren of the Common Life 76.

Catalogue of vices 19, 66, 78, 83, of virtues 19, 83.

Catechetical Sermons 37 ff., 66, 76, 82, 88, 119, 127, 135, 147, and see Luther, Index i.

Catechetics, etymology 3, 451, definition 3, division 7, necessity 6, when first introduced into the university curriculum 138, textbooks 1 f., 151. On the usage of the term compare 3, 75, 87, 117, 451.

Catechism, Luther’s Large, origin 83 ff., date of publication 92, name 92, contents 93, “the best commentary on the Small Catechism” 93, quoted often, esp. in § 26; bibliography 82.

Catechism, Luther’s Small, origin 83 ff., date of publication in tables 91, as book 93, earliest editions 94 ff., textual criticism 81 f., 95, translations 96, 118 f., 184, 330 etc.; evangelical nature 101 ff., its linguistic beauty 109, not polemical 110, not a system 111, 318 f., 344, restricted to the essentials 112, linked with daily life 286, a personal confession 105 f., 286, a summary of the church’s faith 286; even today it is an adequate textbook for religious instruction 311, objections raised against its use 311 ff., faulty methods of treating it 311 f., the proper expository principles 316 ff., hints to a correct interpretation 322–366; the necessity of instruction in the Catechism 286, 309 ff., 498; the correct method of teaching the Catechism 497 ff. and see Method.

Catechism, Notes on parts of Luther’s Small Catechism: Introduction 323, 584 ff., 591 f. Decalogue, Luther’s conception 101 ff., 112, 312 f., 319, 323 ff.; his freedom toward the text 101, 324, 331, 333; his attitude toward the letter of the law 102; one root of morality 102 f., 110, 328; enumeration of the commandments 100, 325; introduction 312 f., 324; the “dass” 330 f.; first commandment 101, 113, 325, 513 ff., 592 ff., second commandment 331, 441, third commandment 100 f., 332 f., fourth commandment 334, fifth commandment 227, 334, 524, sixth commandment 113, 335, seventh commandment 335 f., eighth commandment 336 ff., ninth commandment 313 f., 338 ff., 519 f., tenth commandment 338 ff., conclusion 101, 343 ff., 371; the decalogue as a mirror for sin 107, 112. Creed 103, 314, 346, 605, vital relation to confessor 105, first article 100, 104, 222, 314, 347 ff., 371, 380, 549, second article 104, 110, 315, 358 f., 371, 380, 441, 525, third article 100, 101, 105, 110, 359 ff., 371, 442, 605 ff. Lord’s Prayer 100, 106, 315, 371, introduction 100, 106; second petition 106, fifth petition 100, 106, 441, seventh petition 100, conclusion 100. Baptism
\quellseite{668}
106 f., 110, 316, 372. Confession 107, 372. Lord’s Supper 107 f., 111, 316, 372.

Catechisms: \emph{Medieval} catechism (bibliography 60, 69 ff.,). Weissenburg 60, 66, Freising 60, 66, Alcuin 59, 66, 68, Ratisbon Missionary 63, Notker 60, 66, 68, Mayence 69; Hussite 79, Moravian 69, 73, 79, 86, Raudnitz 79. Waldensian 73, 79. \emph{Evangelical catechisms antedating Luther’s} (bibliography 81): Agricola 89, Althamer 89, 349, Brenz 89, 104, 119, 137, 147, 323, 349, Capito 89, Greater or Lachmann 89, 349, Sam 89. \emph{Evangelical catechisms}: Aquila 353, Baden 147, Bischoff 118, 323, 353, Bonnus 117, Butzer 117, 350, 627, Calvin (French, Latin, German) 122, 157, Caspari 184, Chytraeus 117, Cranmer 353, Danzig 131, Dietrich 127, 190 f., Dresden Cross 136, 190, Dresser, 332, 353, Flensburg 136, Friedrich 118, Gausewitz 190, Frankfort-on-Oder 131, Gesenius 125, 128, Glassius 125, 129, Gotha 134, Hadamarius 353, 365, Han 359, Hanover 141, Heidelberg 116 f., 123, 157 f., (the famous division patterned after a Lutheran model 123), Helmstaedt 182, Hercko 353, Herford 136, Hiemer 137, Huber 332, 350, 499, Jacobs 186, Joachimstal 337, 356, 499, Jud 122, Judex 117, Kantz 332, 350, Krotel 184, 186, Lasco 122, Loehe 147, 188, Loener 499, Lossius 117, 499 f., Ludwig 186, Luehrs 186, Mann 184, 186, Magdeburg 330, 336 f., 352 f., 356, 358, Marbach 353, Marshal’s Primer 85, Mathesius 117 f., 316, 337, Maurer 353, Meckhart 332, 350, Mecklenburg 125, 136, Mecklenburg-Strelitz 147, Meder 358, Melanchthon 117, Menius 330, Moerlein 118, 353, Nuremberg 125, 128, Opitz, H. 332, 352, 356, Opitz, J. 353, 499, Pontoppidan 190, Quedlinburg 131, Quitman 181, Reu 188, 190 etc., Rhegius 117, Rhenish 147, Sauermann 330, 335 ff., Saxony-Altenburg 96, 131, Seiler 141, Seiss 186, Siber 338, 353, 364, Spaeth 186, Spener 132, 135, Stohlmann 128, 186, Strassburg 117, Stump 188, Sverdrup 191, Tetelbach 118, 125, 190, 337, 342, 350, 363, Trabert 188, Trotzendorf 117, 332, Victorius 499, Vischer 117, Walther, G. 499, Walther, M. 125, 128, 186, Watts 164, Westminster, Large 157, Small 157, 159, Wigand 117, Wuerttemberg 137, 147, Zell 117, 349.

Catechist, the 543 ff.; see also Teacher and Example.

Catechumenate, its development, grades, etc. in the early church 13, 27-30, 33 f., 36 ff., 44.

Catechumens, instruction of: its significance, nature, duration etc. 404, 617, 619 f.; material for it 430 f.; arrangement of a course of study 430 f., 540 f.; bibliography 388, 431.

Centrifugal and centripetal nerves 201.

Character 237, 258.

Charts, catechetical 75, 84, 91, 135.

Childhood, first period 249 ff., second period 254 ff.

Children, objections to religious training of 5.

Children’s services 151, 403 and see Sunday School.

\emph{Christenlehre} 147, 150, 182, 188.

Church attendance, children’s 261, 287, 392.

Church, the, and her confessions as study material for Y.P.S. 432, 538.

\quellseite{669}
Church, the, and her problems as study material for Y.P.S. 432, 538. Church Book as material for instruction 287, 374 ff., 534, and see Church Year, Hymns, Liturgical Training.

Church History as material for instruction 34, 64, 150, 194, 276, 288, 383 ff., 428, 537; nature of such instruction 384, 537; aim 384; outline of a course in Church History for Y.P.S. 432.

Church Orders: Nuremberg Church Order (1533) 96, Order for Liegnitz and Brieg (1535) 623, Ziegenhain Order of Church Discipline (1538) 625, Order for Cassel (1539) 626, 628, Pomeranian Liturgy (1544) 622, 630, Wittenberg Reformation (1545) 120, 618, 630, Order for Hesse (1566, 1574) 627, 632, Order for Brunswick-Wolfenbuettel (1569) 629 f., Order for Electoral Saxony (1580) 619.

Church Polity as material for instruction 150, 277, 288, 387, 433.

Church Year as material for instruction 67, 76, 375 ff.; value and necessity of such instruction 375; the church year as governing the selection of stories for parish and Sunday schools 411, 424 f.

Church’s duty to teach, the, 4 ff., 288, 391, 617.

Circumcision 18, 22.

Close of religious instruction, preliminary 615 ff., final 635 f.

Combination Schools 402, 416 f.

Command, futility of, as means of controlling conduct 246, 248.

Communion of children 53, 75, 79, 616; why rejected 616 f.

Communion privilege 616; see also Confirmation.

Comparison as an element of method 448, 484 f., 566 ff., 587 f., 591, 597 f., 605 f.; see also Method.

Competentes 36, 39 f., 47 ff.

Conceit 228.

Concentric circle plan 409; criticized 409 f.

Concept 205 ff.; concept and conception 205.

Concepts, range of children’s, 212, 250, 252, 292, 439, 455 ff., 478 ff.

Conception 205, 215, 257, 440.

Conclusion (logical thinking) 216, 257.

Concrete 210, 250.

Confession as material for instruction 92, 94 ff., 107.

Confirmation 615 ff.; history of confirmation 49, 62, 74, 120, 134 ff., 150, esp. 615 ff.; Luther on confirmation 618, 630; why no special rite in the Lutheran Church at first 618 f., how it developed 620; the elements of confirmation: examination 620 f. (how to conduct it 621), prayer 621 f., the blessing 622, the confirmation address 623; whether a confession of faith should also be part of the confirmation act 623 ff.; why and how the present practise of confirmation should be modified 631 f.; the day for the first communion 633; the age at which children should be confirmed 617 f.; special courses preparing for confirmation 120, 404 f., 430, 617 f.; see also Catechumens’ instruction.

Confirmed, supplementary care for the, 150, 191, 405, 635; bibliography 388 f.; see also Bible Class and Y.P.S.

Congregationalists 158, 172.

Conscience 225, 257.

\quellseite{670}
Consciousness 203, 206.

Construction, method of 471, 475; criticism 475 f.

\emph{Contra Celsum} by Origen 28 f.

Cooperation of teacher and pupil 143, 220, 260, 451.

\emph{Corpora Doctrinae}, Lutheran, 118.

Correlation of Sunday school and Saturday school 403, 429.

Creed, the, as material for instruction 18, 20 f., 28, 37 ff., 41, 56, 63 ff., 66 f., 74 ff., 78 f., 85, 88, 91, 93, 100 f., 103 ff., 109 f., 111 f., 122 f., 155, 158, 160, 275, 314 f., 317 ff., 346 ff.; its vital significance for the confessor 105 f., 357; its division into 12 parts 64, into 3 parts (Luther) 85, 360.

Culture epoch theory 149.

\emph{Daily Vacation Schools} 192, 402, 429 f.

\emph{Datio salis} 41.

Decalogue as material for instruction 29, 31, 37, 50, 55, 69, 75 ff., 83 ff., 88, 91, 93, 99 f., 101 ff., 107, 110 ff., 122 f., 128, 155 f., 158, 160; Objections raised against its use 312 ff.; Luther’s conception of the decalogue 101 ff., 107, 110, 111 f., 113, 312 ff., 319, 323 ff.

\emph{De catechizandis rudibus} by Augustine 32 ff., 63, etc.

Deduction 216.

Delayed impulses 233.

Delineation of characters, see Bible Biography.

Delusion 211.

Desire 233; desire proper 234; scale of desires 234 f.

Development, physical, of the pupil, 242 ff., 249, 254, 263 f.

Development (in method) 151, 441 ff.; the two forms: analytical 442, and synthetic 442, the analytical to be especially cultivated 443; development in instruction in the catechism 529 f., 584 ff., 591, 592 ff., 605 f.

\emph{Dialogi Sacri} by Castellio 121, 130, 137.

\emph{Dicta Probantia} 157, 369 f., 531 f., see also Bible Passages.

\emph{Didache} 14 f., 17 ff., 22 ff., 47.

\emph{Didascalia} 27.

Discipline 254, 259, 462 f., 472, 481, 545 f., esp. § 22.

Distribution of the material 408 ff.

Divine Service as agency of instruction 34, 88, 90, 227, 391 ff.

Doctrinal instruction see Creed.

Domestic Duties, Table of, 92, 94, 108 f., 334.

Domestic catecumenate, see Home.

Doubt 257, 262, 265 f., 269.

Drill 447, 486, 533.

Duration of baptismal preparation 27, 30, 35, 44, 63.

Dutch Reformed 158.

Duty, sense of, 222.

\emph{\emph{Ecloga}} by Theodul 60, 67 f.

Educational agencies of the Church 385 ff.

Eisenach Conference 95.

Eleusinian Cult 39.

\quellseite{671}
Elimination of Luther’s Catechism demanded 141, 148 f., 311 f., see also Small Catechism.

Emotions 218-231; classified as to origin 219, to contents 219; as concomitant phenomena of the intellectual life 219, 229, 243, 245, 247 f., 252, 257 f.; as organs of evaluation 219 ff., esp. 229; as moving the will 439, 546. Intellectual emotions 220 f., esthetic emotions 221 f., moral emotions 222 ff., religious emotions 224 ff., social emotions 226, self-feelings 228 f.

England, religious education in, 154 ff., 162 ff.

English translations of Luther’s Small Catechism 155, 331.

Environment, influence of, 246, 248, see also Home.

\emph{Epideixis} by Irenaeus 20, 26.

Episcopalians 156, 158, 178, 378.

Erotmatic method defined 450; required by 451; when to be employed 452.

Eschatology as material for instruction 17, 21, 34, 63 f.

Esthetic emotions 221 f., 257.

Evaluation, emotions as organs of, 219, 229.

Evaluation of acts 223, 254, 260, 446, 474, 482 f.

Examination of confirmands see Confirmation.

Example of teacher 208, 223 f., 226 f., 235, 239, 246, 251, 269, 401, 444 f., 472, 543 ff.

Exorcism 29, 31, 41, 54.

Explanation of terms 439.

Expository principles for the exposition of Luther’s Small Catechism 316 ff.

\emph{Fairy tales} as material for religious instruction 149, 289, 291, 409, why to be rejected 291 ff.

Faith, congregational life as a life of, 275; as suggesting the material for instruction 285 f.

Family worship 246, 405.

Fasting 18, 24 f., 36.

Fears 245.

Feebleminded 202, 463.

Feelings see Emotions.

Fides Historica 286.

Fighting instinct 259.

Filial congregations, religious instruction in, 405.

Final theme 443.

Foreign Missions as study material for Y.P.S. 433.

Foreign words 440.

Formal aspects of method 450 ff.

Formal steps, the theory of, 151, 448 f.

France, religious instruction in 70.

Freedom of will 237.

Friendships 252, 257, 262, 267.

\emph{Gangs} 257.

Gary plan 194 f., see Week Day Religious Instruction.

\quellseite{672}
Genuet 230.

Generalization as an element of method 448, 467, 485.

Gentiles, baptismal instruction of 17 f.

\emph{Gloria in Excelsis} as material for instruction 66.

Gospel Book by Otfried 60, 67.

Gospel-hymns 368.

\emph{Gratias} see Benedicite.

Gratitude 227.

Greeley plan 193, see Week Day Religious Instruction.

\emph{Habits} 108, 209, 235 f., 238, 245, 249, 252, 258, 261.

Hail Mary see \emph{Ave Maria}.

Halle institutions 137 f.

Hallucination 211.

Hanse, influence of, upon education 76.

\emph{Heliand} 60, 67.

Heroes, the child's interest in, 256, 260.

Herrnhut 386.

\emph{Hippolyti Canones} 27 f.

\emph{Historia Sacra} by Fabricius 121, 130, 137.

History of salvation as material for instruction 149, 298, 415; brief outline 299 ff.

History of the Church as material for instruction, see Church History.

Holland, religious instruction in 158.

Home, the, as an agency of Christian training 5, 55f., 65, 73, 77, 88, 119, 157 f., 227, 245 f., 262, 279, 389 f., 405, 408.

Honesty, sense of 222.

Honor, sense of 228, 258, 264.

Hussite movement 78 f.

Hymns as material for instruction 76, 90, 122, 129, 131, 150, 161, 247, 253, 375 ff., 430, 434; value 375 f., tunes should be practised 376 ff.; which hymns should be memorized 380 f., method of teaching 535, 611 f., bibliography 374.

Hymns for Beginners, Juvenile Hymnaries 161, 181, 184, 185 f., 376, 416.

\emph{Ideal association} 260, 430, 481.

Ideals 267.

Illusion 211.

Illustrations for the Catechism, collections of, 321.

Image 205 ff., see Concept.

Imagination see Phantasy.

Imitation 245 f., 251.

Immanent review 410, 448.

Immersion 18, 23, 47 f., 64, see Baptism.

Imposition of hands 21, 41, 49, 54, 120, see also Confirmation.

Impulse 232 f., 243.

Inclination 234.

Independence of soul 201, 240.

Individuality of the pupil 229, 259.

\quellseite{673}
\emph{Induction} 216.

Infancy 242; early infancy 242 f., later 243 ff.; bibliography 241.

Infant Baptism: duty and justification 51 f., history of its practise in the early church 52 ff., in the Middle Ages 64; ritual of infant baptism in the early church 54, later 64. References to infant baptism by Irenaeus 52, Tertullian 53, Origen 53, Synod of Carthage 53, Cyprian 53, Apostolic Constitution 53, Gregory Nazianzen 53, Augustine 53, Leo the Great 53, Gregory the Great 53. Instruction of those baptized in infancy 54, 64. Infant baptism as necessitating religious instruction 4.

Infant Schools 391.

Inner life of pupil 200 ff.; bibliography 199 f.

Inner Missions as study material for Y.P.S. 433.

Instinct 233.

Intellect 203 ff.; sensation 203, perception 204, intuition 204, concept 205 ff., motion of concepts 206, association 207, reproduction 207 ff., memory 209, phantasy 210, apperception 212, logical thinking 213 f., conception 215 f., judgment 216, conclusion 216; development in early infancy 243, later infancy 244 f., early childhood 250 f., later childhood 255, 260.

Intellectual emotions 220 f.

Intellectualism 275, 438.

Interaction of soul and body 201.

Interdenominationalism in Sunday Schools 169.

Interest 143, 213, 220 f., 230 f., 279 f., 438.

International Sunday School lessons 180, 186, criticism 186 f., 423; see also Sunday School.

International Sunday School Council of Religious Education 176.

Intuition 143, 204 ff., 250, 253, 255, 291, 439, 443, 478 ff., 494 f., 520.

Introduction into the Bible, see Bible.

\emph{Jesus as perfect fulfilment} of moral law 346.

Jewel, Tetelbach's \emph{Golden Jewel} (1568) 118, 125, 190, 337, 342, 350, 363.

Jews, baptismal instruction of 16 f.

Judgment (logical thinking) 216, 257.

Justice, sense of 222.

Juvenile Hymnals, see Hymns.

\emph{Kinderlehre} 122, 128, 181.

King's Book 155.

Klosterneuburg altar 68.

\emph{Kunsthomilie} 443.

\emph{Kunstkatechese} 510, 520.

\emph{Lancasterian system} 165, 172.

Language, development of 243 f.

Large Catechism see Catechism.

Laws concerning Bible reading in American public schools 394 f.

Laws of fusion 209, of learning 259, of memorizing 244, 255, 260, of reproduction of images 208.

\quellseite{674}
Laying on of hands see Imposition.

Leaders of adolescents, qualifications of 268 ff.

Learning, golden period of 259.

Liturgical training 150, 184, 276, 374 ff., 381.

Liturgy as material for instruction 381.

Logical thinking 213, 440, development in early infancy 243, later infancy 245, 247, early childhood 250, 254, later childhood 256 f., 260.

Lord’s Prayer as material for instruction 22, 39, 41, 49, 56, 63, 65 ff., 74 ff., 83 ff., 88, 92, 95, 100, 106, 108, 122 f., 155, 158, 160, 315, 431; see also Catechism.

Lord’s Supper as material for instruction 22, 29, 49, 78 f., 86, 88, 92, 94, 107 f., 111, 127, 316, 431, 372, 617; see also Catechism and Communion.

Love, Twin Command of 37.

Lutheran chorals and tunes 376 ff.

Lutheran customs 381.

Lutheran Normal Schools 189.

Lutheran Sunday Schools 173, 175, 182 ff.

Lutheran Sunday School Union (1830) 183.

Lutherans in America, Religious Education among, 180 ff.

Luther’s Life as study material for Y.P.S. 432.

\emph{Marriage, examination preceding} 121.

Marshall’s Goodly Primer 154, 156.

\emph{Mandata} by Hermas 24.

Massachusetts School Law (1647) 166.

Material for instruction 285 ff.

Maturity, Christian 274, congregational 275 ff., for communion 620 ff.

Maxim 237, 370.

Memory defined and discussed 209 f., logical and mechanical memory 209; gradual development in childhood 252, 255 f., 259 f., how to memorize 244, 255, 260; on the duty of explaining memory material 143, 155, 210, 253, 260, 311 f., 430, 535; objections to the practise of memorizing Bible passages 369, not valid 369 f.; hymns to be memorized 380; memory material 414 ff. On the practise of memorizing compare also 55 f., 67, 78 f., 114 f., 121, 143, 155, 176, 438, 499 ff.

Mendacity 254.

Method, Principles of 437 ff., bibliography 437.

Method of instruction in Biblical History 465 ff.; statement of aim 448, 467 f., 549 ff., 563, 574, preparation 448, 467, 469 ff., 549 ff., 563, 574 f. presentation 448, 467, 471 ff., 552 ff., 564, 575, penetration 448, 455, 465, 467, 482 ff., 558 ff., 566, application 448, 367, 487 ff., 562; as varying according to grades 489 f., or schools 490 ff.; bibliography 465 f.

Method of instruction in the Catechism 498 ff.; history of method 499 ff., during Reformation times 499 ff., the Pietists’ analyzing method 502 ff., the Rationalists’ Socratic method 505 ff., various methods advocated since 1830 509 ff., the scientific catechesis 510 ff.
\quellseite{675}
the proper method based upon the principle of intuition 520 ff., its two modes of procedure 523, which one is to be preferred 526 f., it must pursue practical aims 528 f.; the various steps: statement of aim 529, developing the specific truth from the intuitive material 529 ff., 584 ff., 591 f., 605, comparison 530, 587 f., 591, 597 f., 605 f., valuation 530 f., 588 f., 592, 603 ff., application 531 f., 589 f., 592, 603 ff., drill and review 533, 608 ff.; it should be Christo-centric 534.

Method of instruction in Hymns, Church History, etc. 535 ff.

Method of instruction employed in American Sunday schools 176 ff.

Methodists 167 f., 172, 378.

Miniature painting 109, 204, 253, 479.

Miracle and morality plays 60, 68.

Missions as material for instruction 150, 193, 227, 315, 321, esp. 386 f.

Mithraism 31 f., 39.

Mnemonic memorizing 209 f.

Moods 230.

Motivation 474, 476 f., 558.

Moral consciousness 223.

Moral emotions 222 ff.

Moravian Brethren 78 f., 386.

Motion of concepts 206.

Mystery cults 31, 39 f.

Narration 250 ff., 477 ff.

Narrowness of consciousness 207 f., 250, 262, 424, 468.

Neighbor, love of 227.

Neo-Herbartīans 291.

Nerves 201.

Neurasthenic and neurotic children 202.

\emph{New England Primer} 160.

\emph{Nuremberg Sermons for Children} 95 f., 120, 155, 312, 333, 336, 338, 341, 345, 350, 357 f., in English (1548) 155.

Office of the Keys 95, authorship 96.

Old Testament History as material for instruction 294 f., why rejected by some 149, 294 f., why it is needed 294 f., what stories to select from the Old Testament 422.

Order of Biblical Books, when and how to learn 372.

Orphism 39.

\emph{Parents’ advantages} 390.

Parish School 397 ff., 417 ff.; necessity and blessing 188, 398 ff., 417 f.; not an “exotic” growth 188, 398 f.; curriculum 409 ff.; bibliography 154, 388.

Partial Parish School 400 f., 416 f.

Passion 234.

\emph{Passionalia} 76; Luther’s \emph{Passionale} 97, see Biblical History.

Pastoral nature of catechumen’s instruction 322.

Pedagogics, textbooks on 2, 200, 437.

Penetration as an element of method 440 ff., 455, 465, 467 ff. 482 ff.,
\quellseite{676}
necessity 482, purpose 482 f., nature 483 f., its “goalwardness” 485, form 486, examples 558 ff., 560; helps 483.

Pennsylvania school conditions in colonial times 165 f.

Perception 204, 243, 250, 439.

Pericope as material for instruction 67, 79, 121, 130, 150, 371.

Period, the class period, how it is to be conducted 539 ff., in Sunday school 541.

Persecutions, influence of, on Christian training 27, 33.

Phantasy 210 f.; development in later infancy 244, 247, early childhood 251, late childhood 256, 260, adolescence 266 f.; as the key to the emotions, 211, 252.

Pictures in religious instruction 244, 246, 252, 260.

Placards used in religious instruction 75, 84, 91.

Play 244, 245, 251.

Plays, religious 60, 68.

Polemics out of place in religious instruction of children 110, 445.

Prayer, training in, see Lord’s Prayer.

Prenatal influence 245.

Preparation as an element of method 448, 467, 469 ff., purpose 469 f., form 471, when to be omitted 471, examples 549 ff., 563, 574 f.

Preparation of the catechist 539 ff., 543 ff.

Presbyterians 157 f., 172, 180.

Presentation as an element of method 439, 448, 467, 471 ff., modes of presentation and their evaluation 471 ff., the correct mode is narration 476 ff., various forms of narration 476 ff., narration must be simple 477, intelligible 478, intuitional 478 ff., vivid 481, and must have movement 481; whether stories are to be told in sections 481; examples 552 ff., 564, 575 ff.; helps and bibliography 488 f.

Pride 228.

Primary and secondary material for instruction 288.

Principles for the exposition of Luther’s Small Catechism 118, 147, 190; unfold fulness of verbal contents 316 ff., restricted to essentials 317 f., not to systematize 318, not to add another explanation on the text proper 320; see also Catechism.

Prohibition of image worship 101 f., 155.

Proof texts see Dicta Probantia.

Propensitas ad malum 239.

Propensity 234.

Proselyte baptism, catechism, and catechumenate 17, 19, 22.

Psychology textbooks 199; child psychology texts 241.

Puberty, development of 266 ff.

Pupil, physical welfare of 202.

Public School, religion in 394 ff.; Bible reading 395 f., criticism 395 ff.; State Laws concerning Bible reading 394 f.

\emph{\emph{Quakers}} 168.

Qualifications of teacher 543 ff., good education 543, love for the children 545, believing personality, sanctified by the Spirit 546, living faith in the creative power of God’s Word 548.

\quellseite{677}
Question, nature of 452; attributes 455 ff., perspicuity 456, appropriateness 457; kinds of questions: analyzing 135, auxiliary 454, causal 454, comparison 454, confessional 451, 455, deciding 454, defining 454, developing 454, didactic 451, 454, examination 451, 454, final 454, preparatory 454; recitation 454, relation 454, Socratic 454, supplementary 454; faulty questions 458 ff.; treatment of answers 461.

Questions and Answers for Prospective Communicants 96.

Questions asked by children 245, 250, 257, 261 f.

\emph{Reason} 214.

Reciprocity between body and soul 201.

Redditio Symboli 42.

Reformed Church 122 f., 158, 333.

Religious education in the public school 192.

Religious feelings 224 f., 245, 257.

Repentance of children 262.

Renunciation 23, 29, 41, 48, 54, 61, 63.

Reproduction of concepts 206 ff., in early infancy 243, later infancy 244, early childhood 252, later childhood 255.

Restlessness 244, 249, 259.

Review 209, 425, 448, 466, 493, 580, 608 ff.

Rewards 248.

Rhymed verses as memory aids 372.

Roman Catholic aim of training 78, 275; catechetical work 123 f., catechisms 116, 124; teachers in public schools 400.

\emph{Sacramentaries}, Gelasian 44, Gregorian 64.

\emph{Sacramentum Catechumenorum} 41.

Salvation Army 378.

Saturday School 402, 429, correlation to Sunday School 429.

Saxon Mirror 102.

Schoenewalde visitation records 91.

School Bibles see Bible.

Schools, various types: Cathedral schools 67, Christian day schools 129 f., 139 f., and see Parish School, convent schools 57, 67, daily vacation schools 192, 402, 429 f., German schools for reading and writing 76, 87, 90, Latin schools 76, 87, 117, 119, 121, neighborhood schools 167, parish schools (Charlemagne’s) 67, Saturday schools 402, 429, summer schools 402, 429 and see Daily Vacation Schools, Sunday schools see Sunday School, verger schools 90, 119, village schools 129, week day religious schools see Week Day Religious Instruction.

Schwenckfelders 159.

Scotland, religious education in 156.

Scrutinies 42, 44, 54.

Secondary material for instruction 383 ff.

Selection of stories for curriculum 410, see also Biblical History.

Self-consciousness, continuity of 201, 240, unity of 201, 240; age at which children attain to 245.

\quellseite{678}
Self-deception, danger of, in religious instruction 136.

Self-esteem 228, 258.

Senior department of Sunday School, material for 428 f.

Sensation 203, 439, classes of sensation 204, in infancy 243 f., childhood 250.

Service see Divine Service.

Seventh Day Baptists 159.

Sex, maturity, etc. 258, 263, 265, 268.

Shame 228.

\emph{Sibylline Books} 14.

\emph{Signatio crucis} 41.

Singing as element of instruction 150, 376 ff., 379 f., see Hymns.

Social feelings 226 f., development in later infancy 245, early childhood 252, 254, later childhood 257, 262, adolescence 267.

Socratic method 142, 151, 438, 450, 505 ff.

Soul, existence of 200 f., 240.

Sponsors and sponsorship 54, 65, 72, 156, 382, 635.

Study material employed in American Sunday Schools 176.

Suicide, where to be treated in the catechism 334.

Summer Schools see Daily Vacation Schools.

Sunday School: History of Sunday schools 119, 123, 146 f., 150, 157 f., esp. 162 ff., in America 165 ff., among Lutherans 183 ff.; types of schools 171, departments 418, Sunday school libraries 171, 173; aim 403, need 403, as a missionary agency 403, as a junior congregation 403, 425, as a preparation for church attendance 392 f., its limitations and deficiencies 399, 420, teachers’ meetings 184, 404, 420, teachers’ institutes 420; curriculum for Lutheran Sunday schools 418 ff., 492, Wartburg Lesson Helps see Reu, Index i.

Sunday School Societies: American S. S. Union 169 ff., 183, Boston Society for the Moral and Religious Instruction of the Poor 168, British S. S. Union 164 f., Evangelical Society of Philadelphia 168, Female Union Society in New York 168, First Day or Sunday School Society 167, International S. S. Association 174, 180, 186 f., 423, London S. S. Society 164, Lutheran S. S. Union 183, New York S. S. Union 168, Pennsylvania S. S. Union 168, Philadelphia Sunday and Adult School Union 168 f., World S. S. Association 174.

Sympathy 226.

Synods and Councils: Synod of Carthage (256) 46, 53, Synod of Agde (506) 44, Council of Bracara (610) 44, Fourth Lateran Council (1215) 75, Synod of Cologne (1279) 74, Synod of Liege (1287) 74, Synod of Toledo (1323) 71, Synod of Lavaux (1368) 71, Council of Trent 124, Synod of Dort (1618) 123, 158.

Synods, Lutheran: Augustana 191, General Council 186 ff., General Synod 186, 188, Iowa 189 ff., 317, 376, cf. Reu, Minnesota 189, Missouri 189 ff., 376, North Carolina 181 f., Norwegian Synod 191, Norwegian Lutheran Church 190, 376, Ohio 190, Pennsylvania Ministerium 184, Pennsylvania Synod 182 f., United Lutheran Church 376, Wisconsin 189 f.

Synthetic development 442 f.

\quellseite{679}
\emph{Table of Christian Life} 76; Table of Domestic Duties 92, 94, 108 f., 334; Table of Prayers 92; see also Charts.

Tales see Fairy Tales.

Talmud 14.

Tasks of the Church as material for instruction 150, 276, 288, 386 f., 428, 431, 433; bibliography 383.

Teacher, the 543 ff.; see also Example; Teachers in the early church 24, 34, 37.

Teachers’ meetings 184, 404, 420.

Teaching function of the Church 4 ff., 288, 391.

Temper 230.

Temperaments 229, 259.

Ten Commandments see Decalogue.

Topics for Luther Leagues 432.

\emph{Traditio evangeliorum} 42, orationis dominicae 42, symboli 42.

Translations of Luther’s Small Catechism, English 155, 331, and see Small Catechism.

Truth, sense of 222, 254.

Two Ways 19, 24.

\emph{Uniform lessons} for Sunday Schools 177 ff., 421 ff.

Unitarians 172, 178.

Universalists 168.

Usages of the Church as study material for Y.P.S. 433.

Utilitarianism 223.

\emph{\emph{Verbum abbreviatum}} (i.e. the Creed) 39.

Vacation Schools 192 f., 402, 429 f.

Vanity 228.

\emph{Vernunft and Verstand} 216, 236.

Virtue, Christian 235.

Visitations, the Church Visitations of 1528, 77, Instructions for Visitors 82, 90.

Vividness of narration 481.

Vocabulary of children 439, 455 f., see also Range of Children’s Concepts.

Volitional Life 231 ff., impulse 233, desire 234, inclination 234, propensity 234, passion 234 f., volition proper 235 f.; development in early infancy 243, later infancy 245, early childhood 252, later childhood 258, 260 f.; on training the will 238, 245, 258, 261, see also Habit.

\emph{Waldenses} 78 f., 385.

Wartburg League Topics 432.

Week Day Religious Instruction 119, 147, 193 ff., 401, 417, types of Week Day Schools 195 f.

Wiclifite movement 78.

Will 231 ff., an independent power of the soul 232, see also Volitional Life.

Wittenberg Catechesis (1571) 118.

\quellseite{680}
Work of the Church as material for instruction see Tasks.

Worship as expression of the congregation’s life 276, 287, as an element of religious instruction 150, 184, 276, 374 ff., 381, see also Liturgical Training and Divine Service.

\emph{Yoke of Commandments} 17; Yoke of the Kingdom of God 17.

Young People’s Society 150, 187, 191, 270, 406, 445; study material for Y. P. S. 372, 375, 428 f., 431 ff.; method of study 537 f., 542; bibliography 389; see also Bible Class and Supplementary Care for the Confirmed.
